\begin{abstract}
Chemical translation evaluation must distinguish general quality from retention of technical
terminology. We describe a segment-level methodology combining Google Patents Chemical abstract
pairs with an article-focused JRC-Acquis legal comparison corpus. The pipeline constructs source
snapshots, extracts reference-side terminology using complementary linguistic and learned approaches,
attaches external evidence, and selects exact reference spans through domain-specific large language
model curation. The proposed generation protocol distinguishes a source-only baseline from an
explicitly reference-derived, oracle terminology condition. Evaluation combines corpus lexical overlap,
segment-level learned similarity, and strict and variant-aware term coverage. Intermediate artifacts
and provenance support auditable comparisons. We report descriptive results from ten current model
runs across chemical patents, legal articles, and legal definitions. Mixed rankings conceal
family-level differences, and few-shot effects vary by model and dataset. These aggregate results
do not establish a causal benefit of oracle terminology or chemically validated annotations.
\end{abstract}